\documentclass[12pt]{scrreprt}
\usepackage{arbeitsblatt}
\usepackage{pifont}

\title{Mensch - Maschine}
\ohead{Informatik G11}

\begin{document}
\section*{Anleitung}
\subsection*{Wann endet das Spiel?}
\begin{itemize}
	\item Spieler erreicht die gegenüberliegende Seite 
	\item Spieler hat keine Figuren mehr
	\item Spieler kann nicht mehr ziehen
\end{itemize}

\subsection*{Erlaubte Züge Mensch}
\begin{itemize}
	\item Mensch darf nicht mit linker Figur beginnen -- liefert nur symmetrische Positionen, aber viel mehr Zugkarten!
	\item Geradeaus, wenn frei
	\item Schräg zum Schlagen
\end{itemize}

\subsection*{Erlaubte Züge Maschine}
Vereinfachte Version zum Zeitsparen!

\begin{enumerate}
	\item Suche Situation in der Situationsübersicht
	\item Finde mögliche Züge in der Zugübersicht
	\item Wähle zufälligen Zug (z.B. Verdeckte Farbkarte wählen)
	\item Protokolliere den Zug (z.B. A1 rot)
\end{enumerate}

\subsection*{Ablauf}
\begin{itemize}
	\item Der Mensch beginnt immer
	\item Es wird abwechselnd gezogen, bis der Siger ermittelt ist.
	\item Verliert die Maschine, so wird der letzte Zug (Farbpunkt) durchgestrichen
	\item Der Sieger wird aufgeschrieben.
\end{itemize}

\newpage

\section*{Beispiel-Ablauf}

\begin{itemize}
	\item Mensch rechts $\Rightarrow$ A1
	\item Maschine RGY $\Rightarrow$ R (Figurverlust Mensch)
	\item Mensch mittig $\Rightarrow$ B5
	\item Maschine RGBY $\Rightarrow$ G
	\item Mensch mittig $\Rightarrow$ Sieg
	\item Streichen der Farbe Grün auf B5
	\item Sieger Mensch notieren
	\item \textbf{Nächste Runde}
\end{itemize}


\end{document}
